% ============================================================
% CAPITOLO 25 - QUADRO GENERALE DELLE SETTE SPECIE DI SETTIMA
% ============================================================
\chapter{Quadro generale delle sette specie di settima}
\label{chap:sette_specie}

Dopo la settima di dominante conviene considerare l'intera famiglia delle quadriadi: le diverse combinazioni di terze sovrapposte che nascono sui gradi della scala. Il comportamento delle dissonanze, la loro preparazione e la loro risoluzione sono quelli del Capitolo \ref{chap:teoria_dissonanza}.

\section{Classificazione}
\label{sec:classificazione_specie}

Si distinguono sette specie di accordi di settima, determinate dalla triade di base e dalla qualità della settima rispetto alla fondamentale.

\begin{enumerate}
    \item \textbf{Prima specie:} triade maggiore e settima minore (settima di dominante).
    \item \textbf{Seconda specie:} triade minore e settima minore.
    \item \textbf{Terza specie:} triade diminuita e settima minore (settima semidiminuita).
    \item \textbf{Quarta specie:} triade maggiore e settima maggiore.
    \item \textbf{Quinta specie:} triade diminuita e settima diminuita (settima diminuita).
    \item \textbf{Sesta specie:} triade minore e settima maggiore.
    \item \textbf{Settima specie:} triade aumentata e settima maggiore.
\end{enumerate}

\begin{esempio}[Le sette specie costruite su Do]
	\label{ex:sette_specie}
	Le prime quattro specie e le tre restanti, costruite sulla stessa fondamentale. Sotto ogni accordo sono indicate le qualità delle tre terze sovrapposte (M maggiore, m minore).
	\musica{es_cap13_10}
	\musica{es_cap13_11}
\end{esempio}

In un corso di armonia complementare ci si concentra sulle specie più comuni e si lasciano da parte le più rare, che non trovano applicazione negli esami (Capitolo \ref{chap:specie_rare}).

\section{Collocazione sui gradi delle scale}
\label{sec:collocazione_scale}

Gli accordi di settima si formano sovrapponendo terze \textbf{restando nei suoni della scala}: a ogni grado corrisponde un accordo di specie determinata. Le sigle sotto gli accordi indicano il grado e la specie.

\begin{esempio}[Settime sui gradi di Do maggiore]
	\label{ex:scala_do}
	\musica{es_cap13_12}
\end{esempio}

\begin{esempio}[Settime sui gradi di La minore naturale]
	\label{ex:scala_lan}
	Nel minore naturale la dominante è \emph{minore} (\sigla{V}{7}{} di seconda specie) e manca la sensibile: per questo nella pratica tonale si ricorre alla scala armonica.
	\musica{es_cap13_13}
\end{esempio}

\begin{esempio}[Settime sui gradi di La minore armonica]
	\label{ex:scala_la_arm}
	Sono gli accordi caratteristici del modo minore: \sigla{V}{7}{} (dominante, con la sensibile Sol$\sharp$), \sigla{VII}{7}{} (settima diminuita di sensibile) e \sigla{II}{7}{} (semidiminuita, preparazione alla dominante). Gli accordi con quinta eccedente (\sigla{III}{7}{}) e minore-maggiore (\sigla{I}{7}{}) sono rari.
	\musica{es_cap13_14}
\end{esempio}

\begin{esempio}[Settime sui gradi di La minore melodica ascendente]
	\label{ex:scala_la_mel}
	\musica{es_cap13_15}
\end{esempio}

\begin{center}
\begin{tabular}{ccccc}
\toprule
\textbf{Grado} & \textbf{Do maggiore} & \textbf{La min.\ naturale} & \textbf{La min.\ armonica} & \textbf{La min.\ melodica asc.} \\
\midrule
I & 4\textsuperscript{a} & 2\textsuperscript{a} & 6\textsuperscript{a} & 6\textsuperscript{a} \\
II & 2\textsuperscript{a} & 3\textsuperscript{a} & 3\textsuperscript{a} & 2\textsuperscript{a} \\
III & 2\textsuperscript{a} & 4\textsuperscript{a} & 7\textsuperscript{a} & 7\textsuperscript{a} \\
IV & 4\textsuperscript{a} & 2\textsuperscript{a} & 2\textsuperscript{a} & 1\textsuperscript{a} \\
V & 1\textsuperscript{a} & 2\textsuperscript{a} & 1\textsuperscript{a} & 1\textsuperscript{a} \\
VI & 2\textsuperscript{a} & 4\textsuperscript{a} & 4\textsuperscript{a} & 3\textsuperscript{a} \\
VII & 3\textsuperscript{a} & 1\textsuperscript{a} & 5\textsuperscript{a} & 3\textsuperscript{a} \\
\bottomrule
\end{tabular}
\end{center}

\section{Richiami su prima e terza specie}
\label{sec:richiami_specie}

Le due specie più importanti si trattano in altri capitoli:

\begin{itemize}
    \item \textbf{Prima specie (settima di dominante):} Capitolo \ref{chap:settima_dominante}. Non richiede la preparazione della settima.
    \item \textbf{Terza specie (settima di sensibile, semidiminuita):} Capitolo \ref{chap:vii7_ii7}. Sul VII grado del maggiore la preparazione della settima non è obbligatoria; lo è invece sul II grado del minore. In ogni caso la settima scende di grado. La funzione è di dominante sul VII grado del maggiore, di sottodominante sul II grado del minore.
\end{itemize}
